% =============================================================================
% Patch: §6 Transverse projection + §7 Actual-zeta transverse suppression
%        (referee-revised, with §7 connectivity type-fix)
%
%   sec:transverse-projection
%   sec:actual-zeta-transverse-suppression
%
% Target file:   rh/rh_rebuild.tex
% Action:        this patch is NOT to be merged into rh_rebuild.tex yet.
%                It is staged under rh/patches/ for review.
%
% Scope.
%   §6:
%     - Defines only the two-point scalar transverse projection on M_2(R).
%     - Makes the scalar channel used by §5 and §7 type-safe:
%           A_2(X) := A_toy(Pi_trans X) = A_toy(X) = x_{12} + x_{21}.
%     - Records the quotient-representative identity
%           Pi_trans X = (A_toy(X)/sqrt(2)) e_sym
%       and the HS-norm relation ||Pi_trans X||_HS^2 = A_toy(X)^2 / 2.
%     - Does NOT define the four-point projection Pi_trans^{(4)} or any
%       four-point value subspace; cf. the new
%       rem:transverse-projection-four-point-not-supplied.
%     - Does NOT import v1's M(d), A_val, or Psi_{x,d} calibration
%       machinery (cf. rem:toy-phase-v1-compatibility in §5).
%
%   §7 (connectivity type-fix):
%     - Replaces the ill-typed expression ||Pi_trans A_2(X)|| (with A_2
%       scalar-valued) by the scalar |A_2(X)| = |A_toy(Pi_trans X)| =
%       |A_toy(X)| <= Q^{-B_{z,2}}.
%     - Restates the four-point hypothesis using a separate four-point
%       transverse functional A_{4,trans}, rather than borrowing the
%       two-center scalar projection of §6.
%     - Updates the four-point wip remark to mark Pi_trans^{(4)} and the
%       four-point value subspace as separate inputs not supplied by §6.
%
% Replaces in rh/rh_rebuild.tex:
%   §6 stub (the abstract-conditions definition + rem:wip-transverse-projection),
%   and the two §7 hypothesis statements (hyp:zeta-suppression-two-point,
%   hyp:zeta-suppression-four-point) plus the §7 orientation prose.
%
% Closes:
%   rem:wip-transverse-projection.
%
% Status (when merged):
%   §6:
%     \statusmig      not-started -> migrated
%     \statusreferee  not-started -> in-progress
%     \statussympy    not-started -> verified
%         [rh/sympy/sec-transverse-projection.py]
%     \statussim      not-started -> verified
%         [rh/simulation/sec-transverse-projection.py]
%     \statuslean     not-started -> n/a
%   §7: only orientation+hypotheses are type-fixed; the
%       sec:actual-zeta-transverse-suppression status pills are unchanged
%       since the suppression is still a hypothesis, not a proved theorem.
%
% Verification:
%   - rh/sympy/sec-transverse-projection.py
%       verify_value_subspace_kernel_of_A_toy,
%       verify_explicit_basis_orthonormality,
%       verify_transverse_projection_explicit_formula,
%       verify_pi_trans_idempotent,
%       verify_pi_trans_self_adjoint,
%       verify_pi_trans_kernel_equals_V_val,
%       verify_quotient_representative_and_norm_relation,
%       verify_pi_trans_operator_norm,
%       verify_A_toy_descends_through_pi_trans,
%       verify_A2_type_handoff,
%       verify_A2_quotient_well_defined,
%       verify_A_val_v1_lies_in_V_val.
%   - rh/simulation/sec-transverse-projection.py
%       check_pi_trans_explicit_formula_mp,
%       check_pi_trans_idempotent_mp,
%       check_pi_trans_kernel_on_V_val_basis,
%       check_A_toy_descends_through_pi_trans_mp,
%       check_A2_type_handoff_mp,
%       check_quotient_decomposition_mp,
%       check_pi_trans_operator_norm_unity_mp,
%       check_pi_trans_on_actual_zeta_block,
%       check_pi_trans_on_toy_block_off_line.
% =============================================================================

% =============================================================================
% (1) §6: replace the existing stub (def:transverse-projection abstract
%     conditions + rem:wip-transverse-projection).
% =============================================================================

\section{Transverse projection}
\label{sec:transverse-projection}
\statuspaper{LIVE}\quad
\statusmig{migrated}\quad
\statusreferee{in-progress}\quad
\statussympy[rh/sympy/sec-transverse-projection.py]{verified}\quad
\statussim[rh/simulation/sec-transverse-projection.py]{verified}\quad
\statuslean{not-started}
\par\medskip

The scalar toy anomaly functional \(\mathcal A_\toy\) of
Definition~\ref{def:toy-anomaly-functional} fixes the two-point value
subspace as its kernel.  The transverse projection in this section is
the Hilbert--Schmidt orthogonal projection of the Gram-normalized
\(2\times2\) matrix space onto the one-dimensional complement of that
kernel.  This gives a canonical two-point \(\Pi_\trans\) compatible
with the scalar visibility lower bound of
Theorem~\ref{thm:toy-anomaly-visibility}.  The projection has
Hilbert--Schmidt operator norm one, contributes no polynomial loss,
and allows the actual-zeta two-point suppression statement of
Section~\ref{sec:actual-zeta-transverse-suppression} to be formulated
in the same scalar channel as the toy lower bound.  This section does
not define the four-point projection \(\Pi_\trans^{(4)}\).

\subsection{Value subspace and transverse projection}
\label{subsec:transverse-projection-construction}

We work on \(M_{2}(\mathbb R)\) with the Hilbert--Schmidt inner product
\(
        \langle X,Y\rangle_{\mathrm{HS}}
        :=\operatorname{tr}(X^\top Y).
\)

\begin{definition}[Value subspace]
\label{def:value-subspace}
The two-point scalar value subspace is
\begin{equation}
\label{eq:V-val}
        \mathcal V_\val
        \;:=\;\bigl\{X\in M_{2}(\mathbb R):x_{12}+x_{21}=0\bigr\}
        \;=\;\ker\mathcal A_\toy.
\end{equation}
\end{definition}

\begin{lemma}[Explicit basis for \(\mathcal V_\val\)]
\label{lem:value-subspace-explicit-basis}
\(\mathcal V_\val\) is a three-dimensional subspace of
\(M_2(\mathbb R)\) with orthonormal basis
\begin{equation}
\label{eq:V-val-basis}
        e_{11}=\begin{pmatrix}1&0\\0&0\end{pmatrix},
        \qquad
        e_{22}=\begin{pmatrix}0&0\\0&1\end{pmatrix},
        \qquad
        e_{\mathrm{anti}}=\frac1{\sqrt2}
        \begin{pmatrix}0&1\\-1&0\end{pmatrix}.
\end{equation}
Its one-dimensional orthogonal complement is
\begin{equation}
\label{eq:V-val-perp-basis}
        \mathcal V_\val^\perp=\mathbb R\,e_\sym,
        \qquad
        e_\sym:=\frac1{\sqrt2}
        \begin{pmatrix}0&1\\1&0\end{pmatrix}.
\end{equation}
\end{lemma}

\begin{proof}
The four matrices \(e_{11},e_{22},e_{\mathrm{anti}},e_\sym\) form an
orthonormal basis of \(M_2(\mathbb R)\) by direct calculation.  The
first three have \(x_{12}+x_{21}=0\), so they lie in
\(\mathcal V_\val\).  The matrix \(e_\sym\) has
\(x_{12}+x_{21}=\sqrt2\ne 0\) and is orthogonal to the first three, so
it spans \(\mathcal V_\val^\perp\).
\end{proof}

\begin{remark}[Four-point projection not supplied here]
\label{rem:transverse-projection-four-point-not-supplied}
This section defines the two-center projection on \(M_2(\mathbb R)\).
The four-center projection \(\Pi_\trans^{(4)}\) appearing in
Hypothesis~\ref{hyp:zeta-suppression-four-point} and
Theorem~\ref{thm:four-point-rigidity} is a separate object.  No
four-point quotient, four-point value subspace, or four-point scalar
functional is proved in this section.
\end{remark}

\begin{definition}[Transverse projection]
\label{def:transverse-projection}
The two-point transverse projection
\(\Pi_\trans:M_2(\mathbb R)\to M_2(\mathbb R)\) is the
Hilbert--Schmidt orthogonal projection onto
\(\mathcal V_\val^\perp\).  Equivalently,
\begin{equation}
\label{eq:Pi-trans-explicit}
        \Pi_\trans X
        \;=\;\frac{x_{12}+x_{21}}{2}
        \begin{pmatrix}0&1\\1&0\end{pmatrix}.
\end{equation}
\end{definition}

\begin{definition}[Two-point scalar transverse channel]
\label{def:two-point-scalar-transverse-channel}
For a real two-center whitened block \(X\in M_2(\mathbb R)\), define
\begin{equation}
\label{eq:two-point-scalar-channel}
        \mathcal A_2(X):=\mathcal A_\toy(\Pi_\trans X).
\end{equation}
Lemma~\ref{lem:transverse-projection-A-toy-compatibility} gives
\(\mathcal A_2(X)=\mathcal A_\toy(X)=x_{12}+x_{21}\).  Thus
\(\mathcal A_2\) is a scalar functional on two-point matrix
representatives, or equivalently on the quotient
\(M_2(\mathbb R)/\mathcal V_\val\) after choosing the
Hilbert--Schmidt orthogonal representative \(\Pi_\trans X\).
\end{definition}

\subsection{Properties of \(\Pi_\trans\)}
\label{subsec:transverse-projection-properties}

\begin{lemma}[Projection properties]
\label{lem:transverse-projection-properties}
The operator \(\Pi_\trans\) is idempotent, self-adjoint with respect
to the Hilbert--Schmidt inner product, has kernel exactly
\(\mathcal V_\val\), and, as an operator on
\((M_2(\mathbb R),\|\cdot\|_{\mathrm{HS}})\),
\begin{equation}
\label{eq:Pi-trans-norm}
        \|\Pi_\trans\|_{\op,\mathrm{HS}\to\mathrm{HS}}=1.
\end{equation}
\end{lemma}

\begin{proof}
Formula~\eqref{eq:Pi-trans-explicit} is the same as
\(\Pi_\trans X=\langle X,e_\sym\rangle_{\mathrm{HS}}\,e_\sym\).  Hence
\(\Pi_\trans\) is the Hilbert--Schmidt orthogonal projection onto the
nonzero subspace \(\mathbb R e_\sym\).  Idempotence,
self-adjointness, kernel \(\mathcal V_\val\), and operator norm one
follow immediately.
\end{proof}

\begin{lemma}[Quotient representative and scalar norm]
\label{lem:transverse-projection-quotient-representative}
For every \(X\in M_2(\mathbb R)\),
\begin{equation}
\label{eq:Pi-trans-representative}
        \Pi_\trans X=\frac{\mathcal A_\toy(X)}{\sqrt2}\,e_\sym,
        \qquad X-\Pi_\trans X\in\mathcal V_\val,
\end{equation}
and
\begin{equation}
\label{eq:Pi-trans-HS-norm-scalar}
        \|\Pi_\trans X\|_{\mathrm{HS}}^{\,2}
        =\frac{|\mathcal A_\toy(X)|^{\,2}}{2}.
\end{equation}
\end{lemma}

\begin{proof}
The representative identity follows from
\eqref{eq:Pi-trans-explicit} and
\(e_\sym=2^{-1/2}\bigl(\begin{smallmatrix}0&1\\1&0\end{smallmatrix}\bigr)\).
Applying \(\mathcal A_\toy\) to \(X-\Pi_\trans X\) gives zero by
Lemma~\ref{lem:transverse-projection-A-toy-compatibility}, so the
residual lies in \(\mathcal V_\val\).  The norm identity follows from
\(\|e_\sym\|_{\mathrm{HS}}=1\).
\end{proof}

\begin{lemma}[Scalar-channel compatibility]
\label{lem:transverse-projection-A-toy-compatibility}
For every \(X\in M_2(\mathbb R)\),
\begin{equation}
\label{eq:A-toy-pi-compat}
        \mathcal A_\toy(\Pi_\trans X)\;=\;\mathcal A_\toy(X).
\end{equation}
Consequently the toy visibility lower bound
\eqref{eq:toy-anomaly-visibility-A} is unchanged if the toy whitened
block is first replaced by its projected representative
\(\Pi_\trans\widehat\Om_\toy\).  No matrix-norm lower bound is
asserted; only the scalar channel
\(\mathcal A_2=\mathcal A_\toy\circ\Pi_\trans\) is transported.
\end{lemma}

\begin{proof}
By \eqref{eq:Pi-trans-explicit},
\((\Pi_\trans X)_{12}=(\Pi_\trans X)_{21}=(x_{12}+x_{21})/2\).
Therefore
\[
        \mathcal A_\toy(\Pi_\trans X)
        \;=\;\frac{x_{12}+x_{21}}{2}+\frac{x_{12}+x_{21}}{2}
        \;=\;x_{12}+x_{21}
        \;=\;\mathcal A_\toy(X).
\]
The final assertion follows by applying this equality to
\(X=\widehat\Om_\toy(s;m,u,q_0)\) inside the absolute value in
\eqref{eq:toy-anomaly-visibility-A}.
\end{proof}

\subsection{Scope of the chosen projection}
\label{subsec:transverse-projection-scope}

\begin{remark}[No richer quotient is used in the scalar rebuild]
\label{rem:transverse-projection-no-richer-quotient}
The downstream two-point comparison consumes only the scalar
visibility \(|\mathcal A_2(\widehat\Om_\toy)|\) and the matching
actual-zeta scalar upper bound
\(|\mathcal A_2(\widehat\Om_\z)|\).  At this scalar level, the kernel
of \(\mathcal A_\toy\) is the complete value subspace.  The
orthogonal projection of
Definition~\ref{def:transverse-projection} gives the canonical
Hilbert--Schmidt representative of the quotient map
\(M_2(\mathbb R)\to M_2(\mathbb R)/\mathcal V_\val\), with no
polynomial loss.  Richer matrix-level quotient constructions from the
archived local-projective-rigidity machinery are not imported here;
cf.\ Remark~\ref{rem:toy-phase-v1-compatibility}.
\end{remark}

\begin{remark}[Residual gauge]
\label{rem:transverse-projection-residual-gauge}
The retained chart of Definition~\ref{def:local-phase-chart} pins
\(\Ph_T=\theta\).  Additive constants \(\Ph\mapsto\Ph+c\) leave
\(q,q',q''\), hence \(G_{m,\pm}\), \(N_m\), and \(\widehat\Om_\z\),
unchanged.  Affine-jet shifts \(\Ph\mapsto\Ph+at\) change \(q\) to
\(q+a\) and do not preserve the chart class.  Thus there is no
residual gauge action on the two-point matrix representatives beyond
the trivial additive-constant invariance.
\end{remark}

\begin{lemma}[Sufficient two-point projection input]
\label{lem:transverse-projection-sufficient}
The transverse projection \(\Pi_\trans\) and scalar channel
\(\mathcal A_2\) supply the two-point projection input required by
Section~\ref{sec:projective-rigidity}: the toy visibility lower bound
\eqref{eq:toy-anomaly-visibility-A} transports through
\(\Pi_\trans\) by
Lemma~\ref{lem:transverse-projection-A-toy-compatibility}, and the
Hilbert--Schmidt operator-norm bound
\eqref{eq:Pi-trans-norm} contributes no polynomial-in-\(Q\) loss.
The actual-zeta two-point upper bound must be stated in the same
scalar channel, namely as a bound for
\(|\mathcal A_2(\widehat\Om_\z(s;m))|\).
\end{lemma}

\begin{proof}
Combine Lemmas~\ref{lem:transverse-projection-properties},
\ref{lem:transverse-projection-quotient-representative},
and~\ref{lem:transverse-projection-A-toy-compatibility} with
Theorem~\ref{thm:toy-anomaly-visibility} and
Corollary~\ref{cor:toy-visibility-polynomial-weight}.
\end{proof}

% =============================================================================
% (2) §7: connectivity type-fix.  Replaces the existing orientation prose
%     and the two-point + four-point hypothesis statements (and their wip
%     remarks) of sec:actual-zeta-transverse-suppression.
% =============================================================================

\section{Actual-zeta transverse suppression}
\label{sec:actual-zeta-transverse-suppression}
\statuspaper{LIVE}\quad
\statusmig{not-started}\quad
\statusreferee{not-started}\quad
\statussympy{not-started}\quad
\statussim{not-started}\quad
\statuslean{not-started}
\par\medskip

The actual zeta phase \(\Ph\) defines two whitened blocks --- the
two-center block \(\widehat\Om_\z(s;m)\) and a four-center analogue
\(\widehat\Om_\z^{(4)}\) for the symmetric quartet --- on which the
actual-zeta transverse anomaly is suppressed at a strictly larger
polynomial exponent than the corresponding toy lower bound of
Theorem~\ref{thm:toy-anomaly-visibility}.  In the two-point case, the
transverse scalar channel is the one fixed in
Section~\ref{sec:transverse-projection}:
\[
        \mathcal A_2(X)
        \;:=\;\mathcal A_\toy(\Pi_\trans X)
        \;=\;\mathcal A_\toy(X).
\]
The equality is
Lemma~\ref{lem:transverse-projection-A-toy-compatibility}.  Thus
\(\Pi_\trans\) acts on the block before the scalar functional is
applied; the expression \(\Pi_\trans\mathcal A_2(X)\) is not used.

The two-point and four-point comparisons of
Section~\ref{sec:projective-rigidity} consume different forms of this
suppression, stated separately below.

\subsection{Two-point suppression}
\label{subsec:two-point-suppression}

\begin{hypothesis}[Two-point actual-zeta transverse suppression]
\label{hyp:zeta-suppression-two-point}
For the two-center packet associated with \(\rho\) and a symmetry
partner, after the gauge-fixed normalization of
Definition~\ref{def:local-phase-chart},
\begin{equation}
\label{eq:zeta-suppression-two-point-scalar}
        \bigl|\mathcal A_2\bigl(\widehat\Om_\z(s;m)\bigr)\bigr|
        \;\le\; Q^{-B_{\z,2}}
\end{equation}
for an absolute exponent \(B_{\z,2}>A_\toy\) on the retained
subfamily matched to the toy visibility theorem.  Equivalently, by
Lemma~\ref{lem:transverse-projection-A-toy-compatibility},
\[
        \bigl|\mathcal A_\toy\bigl(\Pi_\trans\widehat\Om_\z(s;m)\bigr)\bigr|
        \;=\;
        \bigl|\mathcal A_\toy\bigl(\widehat\Om_\z(s;m)\bigr)\bigr|
        \;\le\; Q^{-B_{\z,2}}.
\]
\end{hypothesis}

\begin{wip}\label{rem:wip-zeta-suppression-two-point}
The two-point case of the structural gap.  Consumed by
Theorem~\ref{thm:two-point-rigidity}.  Admissible proof routes:
\textup{(i)} a standalone route from the local jet-Gram geometry alone;
\textup{(ii)} an energy-supply route consuming the source bound
of Section~\ref{sec:source-energy-supply} via
Lemma~\ref{lem:source-energy-to-suppression}.  v1 prose to mine: the
odd holomorphic transverse channel and corrected finite-\(s\)
holomorphic whitening sections.
\end{wip}

\subsection{Four-point suppression}
\label{subsec:four-point-suppression}

\begin{hypothesis}[Four-point actual-zeta transverse suppression]
\label{hyp:zeta-suppression-four-point}
For the four-center packet associated with the symmetry quartet
\(\{\rho,\bar\rho,1-\rho,1-\bar\rho\}\), let
\(\Pi_\trans^{(4)}\) be a four-point transverse projection and let
\(\mathcal A_{4,\trans}\) be the corresponding four-point transverse
functional, i.e.\ a scalar or finite-dimensional normed output after
the four-point projection has been applied.  Then
\begin{equation}
\label{eq:zeta-suppression-four-point-transverse}
        \bigl\|\mathcal A_{4,\trans}\bigl(\widehat\Om_\z^{(4)}\bigr)\bigr\|
        \;\le\; Q^{-B_{\z,4}}
\end{equation}
for an absolute exponent \(B_{\z,4}>A_\toy^{(4)}\) on the retained
four-point subfamily.
\end{hypothesis}

\begin{wip}\label{rem:wip-zeta-suppression-four-point}
The four-point case of the structural gap.  Consumed by
Theorem~\ref{thm:four-point-rigidity}.  Same admissible proof routes
as Hypothesis~\ref{hyp:zeta-suppression-two-point}.  The projection
\(\Pi_\trans^{(4)}\), the four-point value subspace, and the
functional \(\mathcal A_{4,\trans}\) are not supplied by
Section~\ref{sec:transverse-projection}; they are separate
four-center inputs.  v1 prose to mine: the local projective rigidity
of the one-pair and multi-pair families section, four-point variant.
\end{wip}

\begin{remark}[Standalone and energy-supply routes]
\label{rem:zeta-suppression-routes}
Each suppression hypothesis is admissible if proved either from the
local jet-Gram geometry alone (the standalone route) or by
transferring the source-energy bound of
Section~\ref{sec:source-energy-supply} via
Lemma~\ref{lem:source-energy-to-suppression} (the energy-supply
route).  The two routes are not exclusive; the rigidity theorems of
Section~\ref{sec:projective-rigidity} indicate which hypothesis is
invoked, not which route proves it.
\end{remark}

% =============================================================================
% End of patch.
% =============================================================================
